\section{Introduction}
\label{sec:intro}

Membership inference asks whether a trained model exposes evidence that a particular document was included in its training set \cite{carlini2021extracting, carlini2023quantifying, shi2023detecting, zhang2025minkpp, wang2025recall}. Prior work shows that fine-tuning can leave readable activation traces and that membership can be inferred from hidden representations, perturbation responses, and low-rank activation structure~\citep{galli2024lumia,ferrando2026memtrace,minder2025narrow,huang2026gdrift,zhang2026gss}. A parallel line of work uses sparse autoencoders (SAEs) to interpret and edit these activations, and a growing set of methods edits, suppresses, or ablates SAE features to reduce undesirable model outputs~\citep{wang2025sspu,frikha2025privacyscalpel,suri2025mitigating,muhamed2025dsg}. Such methods can directly control only the information represented in the learned dictionary and its reconstruction. Recent SAE work cautions against treating that dictionary as a complete account of the activation, since reconstruction error can carry structured components~\citep{engels2024dark} and broader evaluations report limits in reconstruction fidelity, feature completeness, and feature canonicity~\citep{karvonen2025saebench,leask2025sparse}. If the SAE encode--decode path preserves membership information, then feature-level editing acts on a representation that contains the membership signal. If reconstruction weakens the signal while the residual preserves it, then feature-level editing may control targeted outputs without controlling membership inference risk.

In this paper, we propose an activation-level recoverability test for fine-tuning membership in transformers. Our methodology involves three contributions. First, we split mean-pooled residual stream activations into the SAE reconstruction and the reconstruction residual, and train membership detectors on each part, so that membership recoverability can be attributed to either side of the SAE decomposition. Second, we introduce an explicit reconstruction-quality validity gate that decides which model-SAE settings are eligible for quantitative comparison, separating settings where the SAE is reliable enough to interpret from settings where it is not. Third, we treat causal localization of the residual signal as a separate question and use feature-level interventions and direction-aware retraining as falsification tests that scope the residual-gap claim.

The empirical studies show that the SAE reconstruction residual carries more membership signal than the SAE reconstruction across the controlled and broader settings we examine. On Pythia~\citep{biderman2023pythia}, independently trained SAEs reproduce the reconstruction-induced drop, providing evidence against SAE initialization as the sole explanation, and Pythia-12B repeats the residual-over-reconstruction ordering at smaller magnitude. A nine-model sweep with the validity gate active shows the same ordering wherever the SAE reconstruction is reliable enough to interpret, while the magnitude depends on partition and model. Control analyses provide evidence against detector fitting and token-level vocabulary artifacts as sufficient explanations. Removing leading membership directions collapses verbatim extraction on the four models tested at erasure ranks 10 and 20, yet membership detection remains high (Appendix~\ref{app:dd_extraction}). These findings indicate that SAE feature-level interventions should not be treated as privacy guarantees unless the reconstruction residual is explicitly audited, since membership evidence can remain recoverable from the residual even when feature-level interventions suppress outputs that appear privacy-relevant.
